%
% Abstract — revised for clarity and examination readiness
%
\chapter{Abstract}
\begin{SingleSpace}
Urban transportation systems generate heterogeneous traffic and sensor data that must be cleaned, validated, and prepared before analysis. This dissertation presents an intelligent assistant for automating that preparation and extends it into \textbf{Flowmatic}, a deployable platform for batch uploads and smart-city pipelines.

The classical study combines multi-dimensional data quality assessment, automated feature engineering, ensemble anomaly detection, and gradient-boosted event classification. On semi-synthetic Astana traffic data, the full preparation pipeline reaches 96.83\% accuracy (F1 = 0.968) with XGBoost, clearly above raw-data and generic AutoML baselines.

Flowmatic implements the same preparation logic in a NestJS/Vue stack orchestrated with Docker Compose. Batch uploads pass through quality checks and cleaning workers backed by PostgreSQL and S3-compatible storage. A smart-city workbench connects simulated sensor sources, Manual/Auto model routing, and medallion lake export. Eight neural models---covering forecasting, anomaly reconstruction, imputation, graph traffic, and severity classification---are evaluated on held-out temporal test splits with three random seeds; the severity classifier reaches macro-F1 = $0.911 \pm 0.017$. Model weights are published on Hugging Face together with reproducible training scripts.

The dissertation links classical preparation research with a working multimodal platform, providing auditable quality metrics, modality-aware inference, and open experimental artifacts.

\end{SingleSpace}
\keywords{Data Preparation, Intelligent Transportation Systems, Smart Cities, Model Routing, Time Series, Anomaly Detection}
